\section{Methodology}
% ===================================  Overivew of robot ===========================
\subsection{Robot Functionality}

The experimental robot is a custom-built, mobile platform designed for testing navigation strategies. 
It employs two \textbf{Metal DC Geared Motors with Encoders} (SKU: FIT0521) for locomotion, driven by a \textbf{TB6612FNG Dual DC Motor Driver}. 
The system is controlled by an \textbf{Arduino Nano V3.0 (ATmega328P CH340 Version)}, with an 	\textbf{ESP8266} module providing wireless TCP communication with a remote computer.
The robot's motors operate at \textbf{6V} with a \textbf{34:1 metal gearbox} and an integrated quadrature encoder, providing \textbf{341.2 counts per revolution (PPR)} of the output shaft.

\begin{table}[h]
    \centering
    \begin{tabular}{ll}
        \toprule
        \textbf{Specification} & \textbf{Value} \\
        \midrule
        Motor Rated Voltage & 6V \\
        Encoder Rated Voltage & 3.3V / 5V \\
        Gear Reduction Ratio & 34:1 \\
        No-load Speed & 210 RPM @ 0.13A \\
        Maximum Efficiency & 2.0 kg$\cdot$cm @ 170 RPM, 2.0W, 0.6A \\
        Maximum Power Output & 5.2 kg$\cdot$cm @ 110 RPM, 3.1W, 1.1A \\
        Stall Torque & 10 kg$\cdot$cm \\
        Stall Current & 3.2A \\
        Encoder Resolution & 11 counts/rev (motor shaft), 341.2 PPR (output shaft) \\
        Dimensions & 52 mm $\times$ 24.4 mm \\
        Weight & 96g \\
        \bottomrule
    \end{tabular}
    \caption{Motor Specifications}
    \label{tab:motor_specs}
\end{table}



\subsection{Robot Navigation and Control}

The robot operates using two primary movement types:
\begin{itemize}
    \item \textbf{Forward motion}: Both wheels turn at equal speed, maintaining a straight trajectory.
    \item \textbf{In-place rotation}: Wheels rotate in opposite directions, adjusting the robot’s heading without displacement.
\end{itemize}

Navigation is performed by sequentially reaching predefined target points $P_1, P_2, \dots, P_n$. 
Given physical constraints, a threshold distance $d_t$ defines successful target acquisition.

The robot’s state vector is given by:
\begin{equation}
\mathbf{x} = \begin{bmatrix} x \\ y \\ \theta \end{bmatrix}
\end{equation}
where $(x, y)$ represents the robot’s position and $\theta$ is its orientation.

\subsubsection{Embedded Control System}

The microcontroller implements a control system that uses encoder-based feedback and proportional speed adjustments.
Real-time wheel rotation is measured by the quadrature encoders driving interrupts and a PWM-based motor control dynamically adjusts speed to maintain balance.

The controller adjusts motor power to keep the wheels rotating at a target speed, but also to minimise the differnce in total rotation between the wheels.
This means that at a given point of time, if a wheel has traveled a larger rotational angle that the other despite currently having equal speeds the power will be reduced.
Each wheel also stops independently once the total target encoder count has been reached.


% ================================= VLP SYSTEM =========================
\subsection{VLP system overview}
\begin{figure}[h!]
    \includegraphics[scale = 0.3]{{figures/vlp_function}}
    \caption{Overview of VLP system}
    \label{fig:acc_grid}
    \end{figure}

The VLP system consists of four transmitters constructed with consumer grade 10W LED luminaires model REX100DL WHWWDIM. 
The system uses discrete tone multiplexing based on IM/DD.
These luminaires are fixed to an aluminium frame at a 1050 mm height, with fixed x–y locations: TX1 (0, 0), TX2 (0, 780), TX3 (760, 780), and TX4 (760, 0).
For the experiments conducted, square wave frequencies of 15 kHz, 17 kHz, 20 kHz, and 23 kHz were inserted by a modulator driver circuit powered by a 30 V power supply. 
The input of the modulator circuit is connected to a function generator which supplies the required square wave. 


The receiver tag consists of a photodiode, trans-impedance amplifier, and active high pass filter. 
The filter eliminates 100 Hz power-line flicker and the DC manifestation of the ambient light. 
A simple inverting amplifier stage then strengthens the filter output to a level suitable for Analog to Digital Conversion (ADC). 
This signal is then fed to a microcontroller that carries out the ADC. 
The ADC sampling rate is 50 kHz.

A VLP measurment is taken by collecting 512 ADC samples, which takes aproximatly 10.4ms at a 50kHz sample rate. 
The sample is then sent over UART to an ESP8266 where it is sent over WiFi using TCP to python scrip running on a host computer.
Processing of the ADC values is implemented on the host computer.
The signal is centered by subtracting the mean, normalized by dividing by the maximum value, and hanning windowed. 
The FFT of the signal is then taken, and the coresponding LED frequency peaks are extracted by searching for the largest peak within 8 FFT bins of the target frequency.


\textbf{Predicting location from FFT peaks}
A MLP uses a vector of the four RSS LED values as input, and ouputs a 2D x, y position vector. 
The input data was scaled using a standard scaler.
The model architecture comprises 200 input neurons, a second dense layer with 50 neurons, and a final hidden layer wtih 100 neurons. 
Dropout layers with a value of 0.1 are used before and after the final hidden layer.
All dense layers use ReLU as the activation function, L2 regularlisation (value of 0.001) and batch normalisation. 
The output layer has a linear activation function.


\textbf{Automated collection of training data}
To enable accurate collection of a large number of VLP RSS readings at known locations a CNC platform was used.
The platform was a 2-axis belt driven setup using three Nema 17 stepper motors.
An OpenBuilds BlackBox Motion Control System (MCS) with Grbl firmware controled the system and a python scipt running on a host computer communicated commands via uart to the system.
The platform has a span of 1000mm x 900mm

\textbf{VLP fingerprint data collection}
Two seperate VLP datasets were collected, one for training and another for validation.
The training set comprised of 1517 collected points which were equally spaced 25mm apart spanng 1000mm x 900mm.
At each point the VLP measurment was repeated 10 times, the RSS values extraced and then averaged to get the RSS value for that location.

The validation locations comprised 528 points that were generated by a uniform grid spacing of 40mm augmented with random offsets in the x and y components. 
The random offsets were drawn from a uniform distribution between 0 and the 40mm step size apart from on the edges which were clamped to ensure the locaitons remained in bounds of the CNC platform.
The random offsets were used to ensure that the VLP model would not tend to 'clamp' predictions to a set grid.

\begin{figure}[h!]
    \includegraphics[scale = 0.8]{{figures/methodology-CNC_fingerprint}}
    \caption{Location of VLP training and validation points}
    \label{fig:acc_grid}
    \end{figure}

\begin{figure}[h!]
    \includegraphics[scale = 0.6]{{figures/methodology-CNC_3d_rss_plot}}
    \caption{RSS values from training data}
    \label{fig:acc_grid}
    \end{figure}


\textbf{Training the VLP model}
The model was trianed using the training dataset, with early stopping based on performance on the validation set with patence of 30 epochs. 
An Adam optimised was used with an exponetial decay learning rate schedule with inital value of 0.01 and decay rate of 0.995.
The batch size was 32 and mean squared error was used for the loss function.

\textbf{VLP model performance}
The model acheived a average error of 21.0mm with a standard deviation of 10.3mm on the validation dataset.

\begin{figure}[h!]
    \includegraphics[scale = 0.6]{{figures/methodology-CNC_validation_errors}}
    \caption{Distribution of VLP errors}
    \label{fig:acc_grid}
    \end{figure}


%-------------------------- Vive tracking  ------------------------
\subsubsection{Use of HTC vive for groundtruth tracking}

The HTC Vive tracking system was employed to accurately determine the robot's position and heading. 
Its effectiveness in research settings has been well established, with previous work reporting a median tracking accuracy of 4.24~mm. 
The system consists of two fixed ``Lighthouse'' base stations and a mobile Tracker.

The Tracker provides six degrees of freedom (6~DoF) localization---delivering both position (\(X\), \(Y\), \(Z\)) and orientation (\(\gamma\), \(\beta\), \(\alpha\)). 
It connects via Bluetooth to a wireless dongle attached to a host PC running SteamVR (\url{https://store.steampowered.com/steamvr}, accessed on 30 March 2021). 
The Python library \texttt{pyopenvr} supplies bindings for the OpenVR SDK (used by SteamVR), which allows for scripted control of the tracking system. 
Furthermore, the Triad OpenVR wrapper simplifies this interface by providing a single Python function to retrieve the Tracker's current position.

The Tracker's location is reported in a local coordinate system that is centered on the first base station, with its axes aligned to that station's orientation. 
This setup necessitates a transformation from the Tracker's native reference frame to the room's coordinate system. 
To derive this transformation, the Tracker is placed at known positions within the test environment, and corresponding points are recorded in both coordinate frames. 
The resulting mapping from the Vive system to room coordinates is then applied to all subsequent positional data.

The CNC machine was used to accurately position the Vive Tracker when capturing the calibration points, which were set at (900, 0), (0, 0), and (0, 900). 
At each of these positions, 10 readings were taken and averaged to compute the transformation. 
A new transformation was calculated at the beginning of each data collection set to eliminate errors resulting from any movement of the Vive base stations.

After collecting the transformation points, the Vive Tracker was transferred to a platform on the mobile robot at the same height, with its orientation fixed by two mounting pins. 
The Tracker's orientation center is located at the middle of the device, where a mounting pin hole exists. 
In contrast, the center of the mobile robot is defined as the midpoint between its two drive wheels along the axle. 
Since the VLP photodiode receiver is positioned at the robot's center (to support VLP-based estimation), the Vive Tracker was offset by placing it 7.0~mm in front of the robot's center. 
Consequently, the Vive-reported coordinates (\(x_{\text{vive}}\) and \(y_{\text{vive}}\)) must be translated to reflect the robot center's location. 
Using the Vive's yaw measurement---which matches the robot's yaw due to the Tracker's fixed orientation---the mapping is performed as follows:

\[
\text{robot\_x} = \text{vive\_x} + 0.07 \times \sin(\text{vive\_yaw})
\]
\[
\text{robot\_y} = \text{vive\_y} + 0.07 \times \cos(\text{vive\_yaw})
\]






%-------------------------- GENERATED PATHS ------------------------
\subsubsection{Collecting navigation data}
The initial training dataset is created by using the real time vive location readings to navigate. 
Target points were randomly generated which the robot would then navigate to. 
To navigate to a point, the robot will compare its bearing to that of the target, make a heading correction and continue to step in that direction.
This is continued until the robot is within a set threshold from the target point.
Once a target point is reached, the next randomly generated point is navigated to.

Each step of the robot consists of taking a new vive reading to determine the groundtruth locaiton and heading of the robot. 
A VLP reading is then taken, which comprises of three sequentianl 512 adc samples. 
The difference between the current heading provided by the vive and the heading towords the target is then calculated
(math here for the calculation)
If the angle is greated than 20 degrees the robot will then initiate a rotation of the required amount. 
Odometry is used to complete the requried turn, with a mapping between encoder step count and actual heading change estimated from a data collection trial. 
After completing the turn the robot then moves forward a set step size which is calibrated to be 100mm. 

The purpose of the data collection is to build up a database of encoder generated movements, turns and steps, along with a VLP location estimates, and the accurate groundtruth for position and heading.



% %-------------------------- Dataset preparation ------------------------
% \subsubsection{Dataset preparation}
% The paths traced by the simulated robot during a simulation run enable the creation of a dataset suitable for training the navigation model. 
% The targets (or labels) for the dataset are the resulting $x$ and $y$ coordinates, along with the heading $\alpha$ at the end of each movement step (heading correction followed by a step).
% The input data consists of the $x$, $y$ location, and heading $\alpha$ for both the encoder-derived estimates; $x_e, y_e, \theta_e$, and the VLP-derived estimates; $x_{vlp}, y_{vlp}, \theta_{vlp}$.


% \begin{tabular}{l|lll|llllll}
% \toprule
% Point \# & \multicolumn{3}{c}{Targets} & \multicolumn{6}{c}{Input}                                          \\
% \midrule
%      & \boldmath$x$     & \boldmath$y$     & \boldmath$\theta$     & \boldmath$x_e$   & \boldmath$y_e$   & \boldmath$\theta_e$   & \boldmath$x_{vlp}$   & \boldmath$y_{vlp}$   & \boldmath$\theta_{vlp}$   \\
% \midrule
% $1$       &$x^1$      & $y^1$     & $\theta^1$     & $x_e^1$   & $y_e^1$   & $\theta_e^1$   & $x_{vlp}^1$   & $y_{vlp}^1$   & $\theta_{vlp}^1$   \\
% $2$       &$x^2$      & $y^2$     & $\theta^2$     & $x_e^2$   & $y_e^2$   & $\theta_e^2$   & $x_{vlp}^2$   & $y_{vlp}^2$   & $\theta_{vlp}^2$   \\

% $\vdots$ & \multicolumn{3}{c}{$\vdots$} & \multicolumn{6}{c}{$\vdots$}               \\
% $n$       & $x^n$    & $y^n$   & $\theta^n$   & $x_e^n$ & $y_e^n$ & $\theta_e^n$ & $x_{vlp}^n$ & $y_{vlp}^n$ & $\theta_{vlp}^n$ \\
% \bottomrule
% \end{tabular}





% ===================================================================================================================================================
% ================================================= METHODOLOGY END =================================================================================
% ===================================================================================================================================================

